\section{Branch Prediction}

\subsection{Why prediction exists}

In the basic pipeline, every taken branch costs a flush because the CPU guessed the next PC would be \texttt{PC+4}. Loops often branch backward many times, so “predict not taken” performs badly on common code.

Branch prediction tries to guess the next PC early, in the fetch stage, before decode and execute know the truth.

\subsection{What must be predicted}

A predictor needs two answers:

\begin{enumerate}
\item \textbf{Direction}: will the branch/jump be taken?
\item \textbf{Target}: if taken, where should fetch go next?
\end{enumerate}

The CPU only has the current PC during fetch, so the predictor stores past behavior indexed by PC.

\subsection{BHT: branch history table}

\texttt{rtl/branch\_predictor.v} contains a BHT with 2-bit saturating counters:

\begin{lstlisting}[language={}]
00 strong not-taken
01 weak not-taken
10 weak taken
11 strong taken
\end{lstlisting}

The top bit is the prediction. A two-bit counter avoids changing opinion after one unusual branch outcome. This is especially useful for loops: a loop branch may be taken many times and not taken once when the loop exits.

\subsection{BTB: branch target buffer}

The BHT predicts direction, but the CPU also needs a target address. The BTB stores:

\begin{lstlisting}[language={}]
valid bit
tag
target address
\end{lstlisting}

The predictor only predicts taken when the BTB hits and the BHT counter says taken. If there is no BTB target, fetch cannot safely redirect.

\subsection{Training the predictor}

Prediction happens combinationally in IF using \texttt{pc\_if}. Update happens synchronously when a control instruction resolves in EX:

\begin{itemize}
\item taken branches increment the counter toward strongly taken;
\item not-taken branches decrement toward strongly not-taken;
\item jumps are forced to strongly taken;
\item taken control transfers allocate/update the BTB target.
\end{itemize}
\subsection{16.6 Integration in \texttt{cpu\_pipe\_bp.v}}

\texttt{cpu\_pipe\_bp.v} carries prediction metadata through the pipeline:

\begin{lstlisting}[language={}]
ifid_pred_taken / ifid_pred_target
idex_pred_taken / idex_pred_target
\end{lstlisting}

When the instruction reaches EX, the CPU compares predicted behavior against actual behavior:

\begin{lstlisting}[language={}]
direction wrong? actual_taken != predicted_taken
target wrong? actual_taken && predicted_taken && actual_target != predicted_target
\end{lstlisting}

If either is wrong, the pipeline flushes younger instructions and redirects to the correct PC.

\subsection{Performance counters}

The branch-predicting core exposes counters such as:

\begin{itemize}
\item total cycles;
\item number of control transfers;
\item number of taken control transfers;
\item number of mispredictions.
\end{itemize}
These make it possible to compare a predict-not-taken pipeline against the BTB/BHT predictor on the same program.

\bigskip
